%%%
%%% Radical "⼙"
%%%
\section*{Radical 26: ``⼙'' (㔾)}\addcontentsline{toc}{section}{Radical 26: ⼙,㔾}\addcontentsline{loh}{figure}{\#\#\#\# 26: ⼙}

%%%%%%%%%% 卫 %%%%%%%%%%
\subsection*{卫}\addcontentsline{loh}{figure}{卫}

\begin{Entry}{卫}{3}{⼙}
  \begin{Phonetics}{卫}{wei4}
    \definition*{s.}{Wei, um estado da Dinastia Zhou | Sobrenome: Wei}
    \definition{s.}{uma palavra usada no nome do lugar | outro nome para um burro}
    \definition{v.}{defender; guardar; proteger}
  \end{Phonetics}
\end{Entry}

\begin{Entry}{卫生}{3,5}{⼙,⽣}
  \begin{Phonetics}{卫生}{wei4sheng1}[][HSK 3]
    \definition{adj.}{bom para a saúde; higiênico; limpo; capaz de prevenir doenças e benéfico para a saúde}
    \definition{s.}{higiene; saneamento; situação limpa}
  \synonymref{厕所}{ce4suo3}
  \synonymref{干净}{gan1jing4}
  \synonymref{洁净}{jie2jing4}
  \synonymref{清洁}{qing1jie2}
  \antonymref{污秽}{wu1hui4}
  \end{Phonetics}
\end{Entry}

\begin{Entry}{卫生巾}{3,5,3}{⼙,⽣,⼱}
  \begin{Phonetics}{卫生巾}{wei4sheng1jin1}
    \definition{s.}{absorvente higiênico; absorvente feminino; um produto de higiene íntima usado por mulheres durante a menstruação}
  \end{Phonetics}
\end{Entry}

\begin{Entry}{卫生厅}{3,5,4}{⼙,⽣,⼚}
  \begin{Phonetics}{卫生厅}{wei4sheng1ting1}
    \definition*{s.}{Departamento de Saúde (da Província)}
  \end{Phonetics}
\end{Entry}

\begin{Entry}{卫生防疫}{3,5,6,9}{⼙,⽣,⾩,⽧}
  \begin{Phonetics}{卫生防疫}{wei4sheng1fang2yi4}
    \definition{s.}{prevenção contra a epidemia}
  \end{Phonetics}
\end{Entry}

\begin{Entry}{卫生局}{3,5,7}{⼙,⽣,⼫}
  \begin{Phonetics}{卫生局}{wei4sheng1ju2}
    \definition*{s.}{Departamento de Saúde | Escritório de Saúde}
  \end{Phonetics}
\end{Entry}

\begin{Entry}{卫生纸}{3,5,7}{⼙,⽣,⽷}
  \begin{Phonetics}{卫生纸}{wei4sheng1zhi3}
    \definition{s.}{papel higiênico | lenços desinfetantes específicos para mulheres durante o período menstrual}
  \end{Phonetics}
\end{Entry}

\begin{Entry}{卫生间}{3,5,7}{⼙,⽣,⾨}
  \begin{Phonetics}{卫生间}{wei4sheng1jian1}[][HSK 3]
    \definition[间,个]{s.}{banheiro; sanitário; \emph{toilette}; quartos com instalações sanitárias em hotéis ou residências}
  \synonymref{厕所}{ce4suo3}
  \synonymref{更衣室}{geng1yi1shi4}
  \end{Phonetics}
\end{Entry}

\begin{Entry}{卫生套}{3,5,10}{⼙,⽣,⼤}
  \begin{Phonetics}{卫生套}{wei4sheng1 tao4}
    \definition[只]{s.}{preservativo; camisinha}
  \end{Phonetics}
\end{Entry}

\begin{Entry}{卫生部}{3,5,10}{⼙,⽣,⾢}
  \begin{Phonetics}{卫生部}{wei4sheng1bu4}
    \definition*{s.}{Ministério da Saúde | Ministério da Saúde Pública}
  \end{Phonetics}
\end{Entry}

\begin{Entry}{卫生球}{3,5,11}{⼙,⽣,⽟}
  \begin{Phonetics}{卫生球}{wei4sheng1qiu2}
    \definition{s.}{naftalina; bolas de naftalina}
  \end{Phonetics}
\end{Entry}

\begin{Entry}{卫生棉}{3,5,12}{⼙,⽣,⽊}
  \begin{Phonetics}{卫生棉}{wei4sheng1mian2}
    \definition{s.}{absorvente | algodão absorvente esterilizado (usado para curativos ou limpeza de feridas) | absorvente tampão; \emph{tampon}}
  \end{Phonetics}
\end{Entry}

\begin{Entry}{卫生署}{3,5,13}{⼙,⽣,⽹}
  \begin{Phonetics}{卫生署}{wei4sheng1shu3}
    \definition*{s.}{Agência de Saúde (ou Escritório, ou Departamento)}
  \end{Phonetics}
\end{Entry}

\begin{Entry}{卫视}{3,8}{⼙,⾒}
  \begin{Phonetics}{卫视}{wei4shi4}[][HSK 7-9]
    \definition[大,家]{s.}{televisão por satélite; abreviação de 卫星电视}
  \seealsoref{卫星电视}{wei4xing1 dian4shi4}
  \end{Phonetics}
\end{Entry}

\begin{Entry}{卫星}{3,9}{⼙,⽇}
  \begin{Phonetics}{卫星}{wei4xing1}[][HSK 5]
    \definition[个,颗]{s.}{satélite; lua; corpos celestes orbitando planetas | satélite artificial | algo que gira em torno de um centro}
  \end{Phonetics}
\end{Entry}

\begin{Entry}{卫星电视}{3,9,5,8}{⼙,⽇,⽥,⾒}
  \begin{Phonetics}{卫星电视}{wei4xing1 dian4shi4}
    \definition{s.}{TV por satélite; televisão por satélite}
  \synonymref{卫视}{wei4shi4}
  \end{Phonetics}
\end{Entry}

%%%%%%%%%% 印 %%%%%%%%%%
\subsection*{印}\addcontentsline{loh}{figure}{印}

\begin{Entry}{印}{5}{⼙}
  \begin{Phonetics}{印}{yin4}[][HSK 6]
    \definition*{s.}{Sobrenome: Yin}
    \definition[个,枚,道,条]{s.}{selo; lacre; carimbo; estampilha | marca; estampa; impressão}
    \definition{v.}{imprimir; gravar | corresponder; conformar; estar em conformidade com}
  \synonymref{戳}{chuo1}
  \synonymref{合}{he2}
  \synonymref{痕}{hen2}
  \synonymref{迹}{ji4}
  \end{Phonetics}
\end{Entry}

\begin{Entry}{印证}{5,7}{⼙,⾔}
  \begin{Phonetics}{印证}{yin4zheng4}[][HSK 7-9]
    \definition{s.}{prova; evidência; coisas usadas para verificar}
    \definition{v.}{verificar; confirmar; corroborar; comprovar; a prova é consistente com os fatos}
  \synonymref{表明}{biao3ming2}
  \synonymref{见证}{jian4zheng4}
  \synonymref{证明}{zheng4ming2}
  \antonymref{编造}{bian1//zao4}
  \end{Phonetics}
\end{Entry}

\begin{Entry}{印刷}{5,8}{⼙,⼑}
  \begin{Phonetics}{印刷}{yin4shua1}[][HSK 5]
    \definition{v.}{imprimir; imprimir textos, imagens, etc. em papel}
  \synonymref{复印}{fu4yin4}
  \end{Phonetics}
\end{Entry}

\begin{Entry}{印刷术}{5,8,5}{⼙,⼑,⽊}
  \begin{Phonetics}{印刷术}{yin4shua1shu4}[][HSK 7-9]
    \definition{s.}{arte da impressão; impressão}[印刷术改变了世界。===A impressão mudou o mundo.]
  \end{Phonetics}
\end{Entry}

\begin{Entry}{印章}{5,11}{⼙,⾳}
  \begin{Phonetics}{印章}{yin4zhang1}[][HSK 7-9]
    \definition[个]{s.}{selo; carimbo; sinete}
  \end{Phonetics}
\end{Entry}

\begin{Entry}{印象}{5,11}{⼙,⾗}
  \begin{Phonetics}{印象}{yin4xiang4}[][HSK 3]
    \definition[种]{s.}{impressão; marca; ideia; os vestígios deixados por coisas objetivas na mente das pessoas}
  \synonymref{记忆}{ji4yi4}
  \end{Phonetics}
\end{Entry}

%%%%%%%%%% 危 %%%%%%%%%%
\subsection*{危}\addcontentsline{loh}{figure}{危}

\begin{Entry}{危}{6}{⼙}
  \begin{Phonetics}{危}{wei1}
    \definition*{s.}{Wei, a décima segunda das vinte e oito constelações em que a esfera celeste foi dividida, consistindo de três estrelas em forma de triângulo obtuso, uma em Aquário e duas em Pégaso | Wei, uma das mansões lunares | Sobrenome: Wei}
    \definition{adj.}{arriscado; inseguro; perigoso | estar gravemente doente; estar morrendo | alto; íngreme}
    \definition{s.}{perigo | cumeeira (de um telhado)}
    \definition{v.}{pôr em perigo; colocar em perigo; comprometer}
  \antonymref{安}{an1}
  \end{Phonetics}
\end{Entry}

\begin{Entry}{危及}{6,3}{⼙,⼃}
  \begin{Phonetics}{危及}{wei1ji2}[][HSK 7-9]
    \definition{v.}{prejudicar; colocar em perigo; comprometer; ameaçar}
  \end{Phonetics}
\end{Entry}

\begin{Entry}{危机}{6,6}{⼙,⽊}
  \begin{Phonetics}{危机}{wei1ji1}[][HSK 6]
    \definition[个,次]{s.}{crise}
  \seealsoref{危急}{wei1ji2}
  \synonymref{风险}{feng1xian3}
  \synonymref{困境}{kun4jing4}
  \synonymref{难关}{nan2guan1}
  \synonymref{危害}{wei1hai4}
  \synonymref{危难}{wei1nan4}
  \synonymref{危险}{wei1xian3}
  \antonymref{安全}{an1quan2}
  \antonymref{安稳}{an1wen3}
  \antonymref{平安}{ping2'an1}
  \antonymref{平稳}{ping2wen3}
  \end{Phonetics}
\end{Entry}

\begin{Entry}{危急}{6,9}{⼙,⼼}
  \begin{Phonetics}{危急}{wei1ji2}[][HSK 7-9]
    \definition{adj.}{crítico; em perigo iminente; em uma situação desesperadora; perigoso e urgente}
  \synonymref{风险}{feng1xian3}
  \synonymref{紧急}{jin3ji2}
  \synonymref{紧迫}{jin3po4}
  \synonymref{紧张}{jin3zhang1}
  \synonymref{迫切}{po4qie4}
  \synonymref{危害}{wei1hai4}
  \synonymref{危机}{wei1ji1}
  \synonymref{危险}{wei1xian3}
  \synonymref{严重}{yan2zhong4}
  \synonymref{要紧}{yao4jin3}
  \antonymref{安全}{an1quan2}
  \antonymref{安稳}{an1wen3}
  \end{Phonetics}
\end{Entry}

\begin{Entry}{危险}{6,9}{⼙,⾩}
  \begin{Phonetics}{危险}{wei1xian3}[][HSK 3]
    \definition{adj.}{arriscado; perigoso}
  \seealsoref{风险}{feng1xian3}
  \synonymref{惊险}{jing1xian3}
  \synonymref{损害}{sun3hai4}
  \synonymref{危害}{wei1hai4}
  \synonymref{危机}{wei1ji1}
  \synonymref{危急}{wei1ji2}
  \synonymref{隐患}{yin3huan4}
  \antonymref{安全}{an1quan2}
  \antonymref{保险}{bao3//xian3}
  \antonymref{平安}{ping2'an1}
  \antonymref{稳妥}{wen3tuo3}
  \end{Phonetics}
\end{Entry}

\begin{Entry}{危害}{6,10}{⼙,⼧}
  \begin{Phonetics}{危害}{wei1hai4}[][HSK 3]
    \definition[个,种]{s.}{prejuízo; perigo; dano}
    \definition{v.}{destruir; prejudicar; pôr em perigo; pôr em risco}
  \seealsoref{损害}{sun3hai4}
  \synonymref{风险}{feng1xian3}
  \synonymref{迫害}{po4hai4}
  \synonymref{破坏}{po4huai4}
  \synonymref{侵害}{qin1hai4}
  \synonymref{伤害}{shang1hai4}
  \synonymref{危机}{wei1ji1}
  \synonymref{危险}{wei1xian3}
  \synonymref{灾害}{zai1hai4}
  \antonymref{保障}{bao3zhang4}
  \antonymref{救援}{jiu4yuan2}
  \antonymref{维护}{wei2hu4}
  \antonymref{有益}{you3yi4}
  \end{Phonetics}
\end{Entry}

\begin{Entry}{危难}{6,10}{⼙,⾫}
  \begin{Phonetics}{危难}{wei1nan4}
    \definition{s.}{perigo e desastre; calamidade}
  \synonymref{艰难}{jian1nan2}
  \synonymref{困苦}{kun4ku3}
  \synonymref{危急}{wei1ji2}
  \end{Phonetics}
\end{Entry}

%%%%%%%%%% 即 %%%%%%%%%%
\subsection*{即}\addcontentsline{loh}{figure}{即}

\begin{Entry}{即}{7}{⼙}
  \begin{Phonetics}{即}{ji2}[][HSK 7-9]
    \definition{adv.}{no presente; no futuro imediato | prontamente; imediatamente}
    \definition{conj.}{e | mesmo; mesmo que}[她瞟了一眼睡着的孩子，随即匆匆离开了。===Ela olhou para a criança adormecida e então saiu correndo. | 即使下雨我也去。===Eu irei mesmo que chova.]
    \definition{v.}{aproximar"-se; alcançar; estar perto | assumir; ascender a; aceitar; começar a se envolver em | ser motivado pela ocasião | estar perto | Literário: ser; significar}
  \end{Phonetics}
\end{Entry}

\begin{Entry}{即可}{7,5}{⼙,⼝}
  \begin{Phonetics}{即可}{ji2ke3}[][HSK 7-9]
    \definition{adv.}{equivalente a; pode então; pode imediatamente; e isso será suficiente (fazer algo); significa 就可以了}
  \seealsoref{就可以了}{jiu4 ke3yi3 le5}
  \end{Phonetics}
\end{Entry}

\begin{Entry}{即使}{7,8}{⼙,⼈}
  \begin{Phonetics}{即使}{ji2shi3}[][HSK 5]
    \definition{conj.}{mesmo; mesmo que; mesmo se; apesar de; expressando uma concessão hipotética | frequentemente usadas com 也 | usado principalmente na linguagem escrita}
  \seealsoref{即便}{ji2bian4}
  \seealsoref{就是}{jiu4shi4}
  \seealsoref{也}{ye3}
  \synonymref{即或}{ji2 huo4}
  \synonymref{假使}{jia3shi3}
  \synonymref{尽管}{jin3guan3}
  \synonymref{就算}{jiu4suan4}
  \synonymref{哪怕}{na3pa4}
  \synonymref{纵然}{zong4ran2}
  \end{Phonetics}
\end{Entry}

\begin{Entry}{即或}{7,8}{⼙,⼽}
  \begin{Phonetics}{即或}{ji2 huo4}
    \definition{conj.}{Literário: mesmo; mesmo que; embora}
  \end{Phonetics}
\end{Entry}

\begin{Entry}{即若}{7,8}{⼙,⾋}
  \begin{Phonetics}{即若}{ji2 ruo4}
    \definition{conj.}{mesmo; mesmo que; embora; mesmo admitindo que; não obstante}
  \end{Phonetics}
\end{Entry}

\begin{Entry}{即便}{7,9}{⼙,⼈}
  \begin{Phonetics}{即便}{ji2bian4}[][HSK 7-9]
    \definition{conj.}{mesmo; mesmo que; embora; mesmo que isso signifique que, em uma situação hipotética ou extrema, o resultado não mudará | frequentemente usadas com 也 | usado com menos frequência}
  \seealsoref{也}{ye3}
  \end{Phonetics}
\end{Entry}

\begin{Entry}{即便是}{7,9,9}{⼙,⼈,⽇}
  \begin{Phonetics}{即便是}{ji2bian4 shi4}
    \definition{conj.}{mesmo que seja}
  \end{Phonetics}
\end{Entry}

\begin{Entry}{即将}{7,9}{⼙,⼨}
  \begin{Phonetics}{即将}{ji2jiang1}[][HSK 4]
    \definition{adv.}{em breve; estar prestes a; estar a ponto de | indica o futuro próximo; não pode ser seguido por substantivos de tempo específico}
  \seealsoref{将}{jiang1}
  \seealsoref{将要}{jiang1yao4}
  \synonymref{必将}{bi4jiang1}
  \synonymref{就要}{jiu4yao4}
  \synonymref{快要}{kuai4yao4}
  \synonymref{马上}{ma3shang4}
  \antonymref{已经}{yi3jing1}
  \antonymref{已然}{yi3ran2}
  \antonymref{早已}{zao3yi3}
  \end{Phonetics}
\end{Entry}

\begin{Entry}{即是}{7,9}{⼙,⽇}
  \begin{Phonetics}{即是}{ji2 shi4}
    \definition{conj.}{significa; isto é; aquilo é}
  \synonymref{等于}{deng3yu2}
  \synonymref{就是}{jiu4shi4}
  \end{Phonetics}
\end{Entry}

%%%%%%%%%% 却 %%%%%%%%%%
\subsection*{却}\addcontentsline{loh}{figure}{却}

\begin{Entry}{却}{7}{⼙}
  \begin{Phonetics}{却}{que4}[][HSK 4]
    \definition{adv.}{mas; contudo; no entanto; enquanto; indica um ponto de virada}
    \definition{v.}{recuar; retroceder | afastar; repelir; desencorajar | declinar; recusar; rejeitar}
    \definition{v.aux.}{usado depois de certos verbos para indicar a conclusão de uma ação, resultado, equivalente a 去 ou 掉}
  \seealsoref{但是}{dan4shi4}
  \seealsoref{掉}{diao4}
  \seealsoref{去}{qu4}
  \synonymref{辞}{ci2}
  \synonymref{但}{dan4}
  \synonymref{拒}{ju4}
  \synonymref{可是}{ke3shi4}
  \synonymref{推}{tui1}
  \synonymref{退}{tui4}
  \antonymref{进}{jin4}
  \antonymref{受}{shou4}
  \end{Phonetics}
\end{Entry}

\begin{Entry}{却是}{7,9}{⼙,⽇}
  \begin{Phonetics}{却是}{que4shi4}[][HSK 6]
    \definition{conj.}{na verdade; no entanto; o fato é\dots; indica um ponto de virada, contrário às suas expectativas anteriores}
  \end{Phonetics}
\end{Entry}

%%%%%%%%%% 卵 %%%%%%%%%%
\subsection*{卵}\addcontentsline{loh}{figure}{卵}

\begin{Entry}{卵}{7}{⼙}
  \begin{Phonetics}{卵}{luan3}[][HSK 7-9]
    \definition[个,只]{s.}{óvulo; ovo; desova | óvulo fertilizado, ovo (de um inseto); zigoto; em entomologia, isso se refere especificamente a um ovo fertilizado, é o primeiro estágio de desenvolvimento no ciclo de vida de um inseto | Dialeto: testículo ou escroto (de um homem); pênis}
  \seealsoref{卵子}{luan3zi3}
  \end{Phonetics}
\end{Entry}

\begin{Entry}{卵子}{7,3}{⼙,⼦}
  \begin{Phonetics}{卵子}{luan3zi3}
    \definition{s.}{óvulo; ovo}
  \end{Phonetics}
\end{Entry}

%%%%%%%%%% 卷 %%%%%%%%%%
\subsection*{卷}\addcontentsline{loh}{figure}{卷}

\begin{Entry}{卷}{8}{⼙}
  \begin{Phonetics}{卷}{juan3}[][HSK 4]
    \definition{clas.}{usado para pequenas coisas enroladas (maço de papel dinheiro, carretel de filme, etc.) | usado para rolos, carretéis, bobinas, etc.}
    \definition[大,小]{s.}{rolo; carretel; bobina}
    \definition{v.}{enrolar; dobrar algo em um cilindro ou semicírculo | varrer; carregar; levar junto | envolver"-se; participar}
  \antonymref{舒}{shu1}
  \end{Phonetics}
  \begin{Phonetics}{卷}{juan4}[][HSK 4]
    \definition{clas.}{usado para capítulos, seções ou volumes; fascículos}
    \definition[份]{s.}{livro; livros e pinturas que são enrolados para coleção; geralmente se refere a pinturas e caligrafia | papel de exame | arquivo; dossiê}
  \end{Phonetics}
\end{Entry}

\begin{Entry}{卷入}{8,2}{⼙,⼊}
  \begin{Phonetics}{卷入}{juan3ru4}[][HSK 7-9]
    \definition{v.}{ser envolvido; estar envolvido; ser atraído para; estar envolvido em}
  \end{Phonetics}
\end{Entry}

\begin{Entry}{卷子}{8,3}{⼙,⼦}
  \begin{Phonetics}{卷子}{juan3zi5}
    \definition{s.}{rolinho primavera; um tipo de prato de massa, feito amassando em folhas finas, cobrindo um lado com óleo e sal, enrolando"-a e cozinhando"-a no vapor}
  \end{Phonetics}
  \begin{Phonetics}{卷子}{juan4zi5}[][HSK 7-9]
    \definition{s.}{prova; prova de exame; um caderno fino ou uma única folha de papel para anotar as respostas durante as provas}
  \end{Phonetics}
\end{Entry}

%%%%%%%%%% 卸 %%%%%%%%%%
\subsection*{卸}\addcontentsline{loh}{figure}{卸}

\begin{Entry}{卸}{9}{⼙}
  \begin{Phonetics}{卸}{xie4}[][HSK 7-9]
    \definition{v.}{descarregar; deitar | remover; despir | livrar"-se de; esquivar"-se (de sua tarefa, responsabilidade) | remover carga ou frete; descarregar a mercadoria transportada do veículo de transporte | desmontar; desmantelar; remover; despojar; desamarrar os arreios do gado | esquivar"-se; evitar; libertar"-se; esquivar"-se da responsabilidade}
  \antonymref{装}{zhuang1}
  \end{Phonetics}
\end{Entry}

%%%%% EOF %%%%%

